\renewcommand{\Anno}{2024}

\renewcommand{\NumeroProblema}{1}

\renewcommand{\Testo}{
Nel piano cartesiano sia $C$ la circonferenza di centro l'origine e raggio unitario; diremo che un insieme $X$ di punti esterni al cerchio unitario ``illumina'' $C$ se, immaginando di posizionare in ogni punto di $X$ una lampadina accesa, tutta la circonferenza $C$ risulta illuminata.

Per ogni numero reale $h>0$ indichiamo con $N(h)$ il minimo numero di punti, tutti a distanza $h$ dalla circonferenza $C$ ed esterni al cerchio unitario, che riescono ad illuminare $C$.

\begin{enumerate}
\renewcommand{\labelenumi}{(\alph{enumi})}

\item  Determinare tutti i valori di $h$ per cui $N(h)=3$

\item  Calcolare $N(1/2)$

\item  Calcolare il limite $$\lim_{h\to0}h\cdot[N(h)]^2$$

\end{enumerate}
}

\renewcommand{\Risposta}{}

\renewcommand{\Aiutino}{}

\renewcommand{\SoluzioneUfficiale}{}

\renewcommand{\SoluzioniAlternative}{}

\renewcommand{\Osservazioni}{}

\renewcommand{\Commenti}{}
